\documentclass[conference]{IEEEtran}

% ─── Packages ────────────────────────────────────────────────────────────────
\usepackage[T1]{fontenc}
\usepackage[utf8]{inputenc}
\usepackage{amsmath,amssymb}
\usepackage{graphicx}
\usepackage{booktabs}
\usepackage{multirow}
\usepackage{array}
\usepackage{xcolor}
\usepackage{hyperref}
\usepackage{cite}
\usepackage{caption}
\usepackage{subcaption}
\usepackage{microtype}

\hypersetup{
  colorlinks=true,
  linkcolor=blue!70!black,
  citecolor=green!50!black,
  urlcolor=blue!60!black,
}

% ─── Title ────────────────────────────────────────────────────────────────────
\title{Interpretability Without Accuracy Cost: Concept Bottlenecks and
Differentiable Rules for Open-Set IoT NIDS}

\author{%
  \IEEEauthorblockN{Ngari Crisphine Macharia}
  \IEEEauthorblockA{Department of Electrical and Computer Engineering\\
  Southern Illinois University Carbondale\\
  c.ngari@siu.edu}
  \and
  \IEEEauthorblockN{Ning Weng}
  \IEEEauthorblockA{Department of Electrical and Computer Engineering\\
  Southern Illinois University Carbondale\\
  nweng@siu.edu}
}

% ─────────────────────────────────────────────────────────────────────────────
\begin{document}

\maketitle

% ─── Abstract ─────────────────────────────────────────────────────────────────
\begin{abstract}
Machine learning (ML) based network intrusion detection systems (NIDS) face
two persistent deployment challenges: lack of interpretability, and poor
handling of novel attack families unseen at training time (the open-set
problem). We present a joint evaluation of two complementary interpretable
architectures applied to IoT traffic classification and open-set detection.
The first, Concept Bottleneck Models (CBMs), routes predictions through a
layer of human-defined traffic concepts and supports test-time intervention.
The second, NeSy-NIDS, encodes domain-expert attack signatures as
differentiable threshold rules with learnable parameters, yielding an
exactly binary rule-activation space via $k$-annealing and the
Straight-Through Estimator. Both architectures are coupled with a
neural fallback gated by a learnable mixing coefficient $\alpha$, and both
use Mahalanobis distance in their respective representation spaces for
open-set scoring. On CTU-IoT-23 and CIC-IoT-2023, neither method incurs
any classification accuracy cost (F1 within $\pm0.0015$ of an unconstrained
MLP). On CTU-IoT-23, rule-space Mahalanobis achieves OOD AUROC\,=\,0.909
($\sigma$\,=\,0.013), competitive with the best CBM variant (0.918). On
CIC-IoT-2023, both bottleneck architectures fall below the MLP baseline
(0.614/0.591 vs.\ 0.858), a result we attribute to a \emph{compression-structural}
effect: reducing 39 raw features to 8 rule or concept dimensions discards
information necessary to separate the 29 unknown attack families in that
dataset, as confirmed by a Decision Tree baseline achieving AUROC\,=\,0.809
on raw features. NeSy-NIDS additionally provides semantically consistent
threshold drift and a controllable interpretability-vs-OOD trade-off via
$\alpha$-regularisation, at zero F1 cost.
\end{abstract}

\begin{IEEEkeywords}
IoT network intrusion detection, concept bottleneck models, differentiable
rule learning, open-set detection, out-of-distribution detection,
interpretable machine learning
\end{IEEEkeywords}

% ─── 1. Introduction ──────────────────────────────────────────────────────────
\section{Introduction}

Internet of Things (IoT) devices now constitute a dominant fraction of
network endpoints. Their minimal patch cycles and heterogeneous firmware
make them attractive targets for botnet recruitment, distributed
denial-of-service (DDoS) amplification, and command-and-control (C\&C)
infrastructure. Network-level intrusion detection systems (NIDS) are a
practical first line of defence: passive on traffic flows, requiring no
modification to end-devices.

Deep learning (DL)-based NIDS achieve strong classification accuracy on
known attacks, but two structural limitations impede operational deployment.
First, the \emph{interpretability gap}: DL models provide no
human-auditable justification for a flagged alert. Without knowing
\emph{which} traffic features triggered a decision, operators cannot
validate alerts, tune policies, or satisfy regulatory requirements~\cite{rudin2019stop,arp2022dos}.
Post-hoc explainers address this symptom but operate outside the model;
recent work has shown that even the best post-hoc NIDS explanation methods
suffer from fidelity, stability, and completeness limitations~\cite{wei2023xnids}.

Second, the \emph{open-set problem}: NIDS trained on a closed set of known
attacks will inevitably encounter novel families at deployment~\cite{arp2022dos}.
A model that confidently assigns novel traffic to a known class provides false
assurance. Existing approaches to this problem either retrain periodically
or apply separate anomaly detectors, neither of which leverages the
model's internal representation for principled open-set scoring~\cite{han2023owad}.

We address both limitations by embedding interpretability \emph{into} the
model architecture and evaluating its effect on open-set detection.
Two complementary architectures are studied:
(i)~\textbf{CBMs}~\cite{koh2020concept}: predictions are forced through
a bottleneck of human-defined traffic concepts, making each intermediate
decision directly auditable and supporting test-time concept intervention;
(ii)~\textbf{NeSy-NIDS}: domain-expert attack signatures are encoded as
parameterised threshold rules trained end-to-end via $k$-annealing and the
Straight-Through Estimator (STE)~\cite{bengio2013ste,liu2023reinmax},
producing an exactly binary rule-activation vector at inference.
Both models are evaluated on CTU-IoT-23~\cite{parmisano2020ctu} and
CIC-IoT-2023~\cite{neto2023ciciot}, two public labelled IoT traffic datasets
spanning 4 and 5 known attack families and 9 and 29 unknown families
respectively.

Distinct from xNIDS~\cite{wei2023xnids}, which generates post-hoc
explanations from a \emph{black-box} DL-NIDS, our architectures are
\emph{inherently} interpretable: the representation space itself is a set
of auditable concepts or binary rules. Distinct from OWAD~\cite{han2023owad},
which detects normality shift (distribution drift over time) in deployed
detectors, we address the complementary problem of identifying
\emph{unknown attack classes} at inference time in a model whose
representation space is human-interpretable.

\smallskip
\noindent\textbf{Contributions:}
\begin{enumerate}
  \item First joint evaluation of CBMs and differentiable rule learning for
    IoT NIDS, demonstrating that neither architecture sacrifices
    classification accuracy versus an unconstrained MLP.
  \item NeSy-NIDS with $k$-annealing and STE achieves 100\% rule crispness
    (all activations binary at inference), enabling direct rule-level audit
    without quantisation overhead.
  \item Discovery and structural explanation of a \emph{compression-induced
    OOD failure} on CIC-IoT-2023: bottleneck compression to 8 dimensions
    discards OOD-discriminative signal present in the 39-dimensional raw
    feature space, confirmed by a Decision Tree baseline.
  \item A controllable $\alpha$-gate providing a deployable
    interpretability-vs-OOD Pareto front, with the operating point
    selectable at inference without retraining.
  \item Concept leakage quantification and elimination via a
    $\gamma$-fidelity regulariser in CBMs, with evidence that reduced
    leakage improves OOD AUROC on CTU-IoT-23.
\end{enumerate}

% ─── 2. Related Work ──────────────────────────────────────────────────────────
\section{Related Work}

\subsection{Interpretability in DL-based NIDS}

Post-hoc methods such as feature attribution~\cite{wei2023xnids} and
concept drift explanation~\cite{yang2021cade} improve operator understanding
of existing DL-NIDS but do not constrain model behaviour.
Wei et al.~\cite{wei2023xnids} (xNIDS) show that sparse group lasso over
network feature groups outperforms LIME and SHAP on fidelity, stability,
and completeness, and translate explanations into actionable firewall rules.
Han et al.~\cite{han2021deepaid} (DeepAID) apply constrained optimisation
to explain unsupervised DL anomaly detectors in security.
Han et al.~\cite{han2024gead} (GEAD) extract global rules from DL anomaly
detectors via regression trees.
Critically, all of these are \emph{post-hoc}: the underlying model remains
opaque. Our work differs by making the model itself interpretable from the
ground up, so that explanations are guaranteed faithful by construction.

\subsection{Concept Bottleneck Models}

Koh et al.~\cite{koh2020concept} introduced CBMs as a two-stage architecture
in which predictions are forced through an intermediate layer of
human-defined concepts, enabling test-time intervention.
Zarlenga et al.~\cite{zarlenga2022concept} propose Concept Embedding
Models (CEMs) that extend CBMs to higher-dimensional concept representations
without sacrificing intervention capability.
Yuksekgonul et al.~\cite{yuksekgonul2023posthoc} show that any pre-trained
network can be retroactively converted to a CBM using post-hoc concept
alignment.
No prior work applies CBMs in the NIDS domain; our paper provides the first
such analysis, including the first quantitative study of concept leakage
in network traffic classification.

\subsection{Differentiable and Neuro-Symbolic Rule Learning}

Petersen et al.~\cite{petersen2022logic} train deep networks of binary
logic gates end-to-end via continuous relaxation and discretisation,
demonstrating that binary differentiable architectures match accuracy of
standard neural networks.
Yang et al.~\cite{yang2022ste} show that the STE can enforce discrete logical
constraints in neural training, providing gradient-based learning of
Boolean conditions.
Pryor et al.~\cite{pryor2023neupsl} (NeuPSL) jointly learn neural and
probabilistic soft logic parameters.
Cheng et al.~\cite{cheng2023ncrl} learn compositional logical rules for
knowledge graph reasoning.
NeSy-NIDS adapts this family of approaches to the operationally specific
domain of IoT network traffic: rules are conjunctions of threshold
conditions over flow features with learnable boundaries, targeting
the interpretability requirements of security analysts.

\subsection{Open-Set and OOD Detection in Security}

Lee et al.~\cite{lee2018simple} establish Mahalanobis distance in the
penultimate feature space as a principled out-of-distribution (OOD)
detector outperforming the maximum-softmax baseline~\cite{hendrycks2017baseline}.
Fort et al.~\cite{fort2021exploring} show that representation quality
dominates OOD detection performance; stronger encoders make Mahalanobis
distance dramatically more effective.
Han et al.~\cite{han2023owad} (OWAD) address a related but distinct
problem---detecting normality shift (gradual distribution drift) in
already-deployed NIDS---via hypothesis testing on output distributions.
Yang et al.~\cite{yang2021cade} (CADE) explain concept drift in security
ML via contrastive analysis.
Yang et al.~\cite{yang2025openset} apply lightweight open-set detection
directly to Industrial IoT (IIoT) intrusion detection.
Our work differs from OWAD and CADE by addressing unknown class detection
at inference time in \emph{inherently interpretable} models, and by
providing an analysis of how representation-space compression affects
OOD detectability.

\subsection{ML Pitfalls and IoT Datasets}

Arp et al.~\cite{arp2022dos} document ten pitfalls of deploying ML in
computer security, including unrealistic closed-world evaluation and
insufficient interpretability---directly motivating our open-set evaluation
and interpretable architecture design.
De Keersmaeker et al.~\cite{dekeersmaeker2023survey} survey 70+ public IoT
network traffic datasets, situating CTU-IoT-23 and CIC-IoT-2023 as two of
the most complete for ML-based NIDS research.
Alotaibi and Maffeis~\cite{alotaibi2024mateen} (Mateen) demonstrate
adaptive ensemble learning for network anomaly detection, building on the
autoencoder-ensemble paradigm and providing a 2024 state of the art for
online network anomaly detection.

% ─── 3. Methodology ───────────────────────────────────────────────────────────
\section{Methodology}

\begin{figure*}[!t]
  \centering
  \includegraphics[width=\textwidth]{figures/architecture.pdf}
  \caption{Architecture overview. \textbf{(A) Concept Bottleneck Model (CBM):}
    Input features are mapped to $K{=}10$ concept scores $\hat{\mathbf{c}}$
    by a 2-layer MLP encoder; a linear head produces class logits; a
    $\gamma$-fidelity regulariser suppresses concept leakage; Mahalanobis
    distance in concept space provides OOD scores.
    \textbf{(B) NeSy-NIDS:} $R{=}8$ rule templates with learnable thresholds
    $\theta_i$ are scored differentiably via a product of sigmoids; $k$-annealing
    ($k\!:\!1\!\to\!10$) and a Straight-Through Estimator (STE) yield exact binary
    activations $\mathbf{a}\!\in\!\{0,1\}^{8}$ at inference; an $\alpha$-gate
    mixes rule logits with an MLP fallback; Mahalanobis distance in rule-activation
    space provides OOD scores.}
  \label{fig:arch}
\end{figure*}

\subsection{Datasets}

\textbf{CTU-IoT-23}~\cite{parmisano2020ctu} provides labelled network flows
from IoT devices in a controlled lab environment.
We use a 37-feature pre-processed version (packet counts, byte counts,
inter-arrival times, flag counts).
Four classes are treated as \emph{known}: Benign, C\&C-HeartBeat, DDoS, and
Okiru; nine attack families are held as \emph{unknown} for OOD evaluation only.

\textbf{CIC-IoT-2023}~\cite{neto2023ciciot} covers 33 attack scenarios
across 105 IoT devices~\cite{dekeersmaeker2023survey}.
After stratified subsampling of the $\approx$40M-row test partition, the
feature set is 39 z-score normalised columns.
Five classes are known; 29 form the unknown set.
Data splits follow a 70/15/15 stratified partition for both datasets.

\subsection{Concept Bottleneck Models}
\label{sec:cbm}

A CBM~\cite{koh2020concept} factors the classifier $f{:}\,\mathbf{x}{\mapsto}y$
into two stages:
\begin{equation}
  \hat{\mathbf{c}} = g(\mathbf{x}),\qquad \hat{y} = h(\hat{\mathbf{c}}),
\end{equation}
where $g{:}\,\mathbb{R}^d{\to}[0,1]^K$ predicts $K$ binary concepts and
$h{:}\,[0,1]^K{\to}\mathbb{R}^C$ maps concepts to class logits.
We define $K{=}10$ domain-derived flow concepts (e.g., \textit{is\_high\_rate},
\textit{has\_beaconing\_iat}, \textit{has\_incomplete\_handshake}), approximated
as soft binary targets from threshold rules on raw features.

\textbf{Training variants.}
\emph{JointCBM} ($\gamma{=}0$): standard joint training.
\emph{JointCBM} ($\gamma{=}0.5$): joint training with concept-fidelity
regularisation:
\begin{equation}
  \mathcal{L} = \mathcal{L}_{\text{CE}} + \lambda\,\mathcal{L}_{\text{BCE}}
    + \gamma\,\|\hat{\mathbf{c}} - \mathbf{c}_{\text{gt}}\|^2,
  \label{eq:cbm_loss}
\end{equation}
which penalises concept activations that deviate from their ground-truth
binary targets, mitigating \emph{concept leakage}---the encoding of
class-discriminative information beyond the specified concepts.
\emph{SequentialCBM}: $g$ trained first and frozen, then $h$ trained
on $\hat{\mathbf{c}}$.

\textbf{OOD scoring.}
Concept activation vectors from training samples are used to fit
per-class Gaussians $\mathcal{N}(\boldsymbol{\mu}_c,\boldsymbol{\Sigma})$
with shared regularised covariance (Mahalanobis
detector~\cite{lee2018simple,fort2021exploring}):
\begin{equation}
  s(\mathbf{x}) = \min_c\,(\hat{\mathbf{c}}-\boldsymbol{\mu}_c)^\top
  \boldsymbol{\Sigma}^{-1}(\hat{\mathbf{c}}-\boldsymbol{\mu}_c).
  \label{eq:mahal}
\end{equation}

\subsection{NeSy-NIDS: Differentiable Rule Learning}
\label{sec:nesy}

\subsubsection{Rule Templates}

$R$ rule templates are defined a priori from domain knowledge of IoT attack
traffic. Each template specifies a set of feature indices, condition
directions ($>$ or $<$), and initial threshold values.
Eight rules are defined per dataset. Selected CTU-IoT-23 rules are listed
in Table~\ref{tab:rules_ctu}; CIC-IoT-2023 rules in Table~\ref{tab:rules_cic}.

\begin{table*}[!t]
\renewcommand{\arraystretch}{1.1}
\caption{NeSy-NIDS Rule Templates for CTU-IoT-23 (left) and CIC-IoT-2023 (right, features z-scored).
  Initial thresholds are derived from domain knowledge; $\theta$ values are updated during training.}
\label{tab:rules_ctu}
\label{tab:rules_cic}
\centering
\footnotesize
\begin{tabular}{llll@{\hspace{1.5em}}llll}
\toprule
\multicolumn{4}{c}{\textbf{CTU-IoT-23}} &
\multicolumn{4}{c}{\textbf{CIC-IoT-2023}} \\
\cmidrule(r){1-4}\cmidrule(l){5-8}
\textbf{Rule} & \textbf{Conditions} & \textbf{Target} & &
\textbf{Rule} & \textbf{Conditions} & \textbf{Target} & \\
\midrule
\texttt{ddos\_high\_rate}    & Rate$>$1000            & DDoS    &&
\texttt{syn\_flood\_ddos}    & rate$>$0.15, syn$>$1.5 & DDoS-SYN \\
\texttt{cc\_hb\_beacon}      & $1.5{<}$IAT$<5$, $0.1{<}$dur$<5$ & C\&C-HB &&
\texttt{udp\_mirai\_flood}   & udp$<{-}1.0$, rate$>$0 & Mirai \\
\texttt{incompl\_handshake}  & syn$>$2, ack$<$1       & C\&C-HB &&
\texttt{recon\_port\_scan}   & fwd\_iat$>$1.5         & Recon \\
\texttt{okiru\_low\_rate}    & Rate$<$500, IAT$>$2.5  & Okiru   &&
\texttt{vuln\_scan}          & pkt$>$1.0, rate$>$0    & Scan \\
\texttt{asymmetric\_flood}   & src$<$1, dst$>$1       & DDoS    &&
\texttt{benign\_low\_rate}   & rate$<$0, syn$<$0      & Benign \\
\texttt{persistent\_conn}    & IAT$>$2.5              & Benign  &&
\texttt{large\_pld\_flood}   & pld$>$0.4, rate$>$0    & Flood \\
\texttt{benign\_probe}       & pkt$<$2, bytes$<$100   & Benign  &&
\texttt{hi\_entropy\_scan}   & ent$>$0.5, iat$>$0.5   & Scan \\
\texttt{ddos\_short\_burst}  & Rate$>$1000, IAT$<$0.1 & DDoS    &&
\texttt{short\_conn\_flood}  & dur$<{-}0.003$, rate$>$0.1 & Flood \\
\bottomrule
\end{tabular}
\end{table*}

\subsubsection{Differentiable Rule Score}

Given input $\mathbf{x}$, the soft activation of rule $r$ at steepness $k$
is:
\begin{equation}
  a_r(\mathbf{x};k) =
    \prod_{i\in\mathcal{I}_r^{>}}\!\sigma\!\bigl(k(x_i{-}\theta_i)\bigr)\cdot
    \prod_{j\in\mathcal{I}_r^{<}}\!\sigma\!\bigl(k(\theta_j{-}x_j)\bigr),
  \label{eq:rule_score}
\end{equation}
where $\sigma$ is the logistic sigmoid and $\theta_i$ are \emph{learnable}
thresholds. The product-of-sigmoids formulation implements a differentiable
conjunction (AND) of threshold conditions~\cite{petersen2022logic,yang2022ste}.

\subsubsection{$k$-Annealing and Straight-Through Estimator}

$k$ is annealed linearly from $k_0{=}1$ to $k_T{=}10$ over the first 30 epochs:
\begin{equation}
  k(t) = k_0 + \min\!\bigl(\tfrac{t}{T_{\mathrm{warm}}},1\bigr)(k_T{-}k_0).
\end{equation}
At $k{=}10$ the sigmoid approximates a step function, but gradients remain
nonzero.
To enforce exact binary activations at inference, we apply the
STE~\cite{bengio2013ste,liu2023reinmax}:
\begin{equation}
  a_r^{\mathrm{hard}} = a_r + \bigl(\mathbf{1}[a_r{>}0.5] - a_r\bigr)
    \big|_{\mathrm{detach}},
  \label{eq:ste}
\end{equation}
passing the rounded value forward while routing gradients through the
continuous activation. This achieves 100\% \emph{rule crispness}---the
fraction of activations in $[0,0.1)\cup(0.9,1]$.

\subsubsection{Architecture and $\alpha$-Gate}

A RuleBank layer stacks all $R$ rule activations into
$\mathbf{a}(\mathbf{x})\in\{0,1\}^R$ (hard) or $[0,1]^R$ (soft, training).
A linear head maps activations to rule logits
$\mathbf{l}_{\mathrm{rule}}{=}\mathbf{W}\mathbf{a}{+}\mathbf{b}$.
A parallel two-layer multi-layer perceptron (MLP) on raw features
produces neural logits $\mathbf{l}_{\mathrm{neural}}$.
A scalar mixing gate $\alpha{=}\sigma(w_\alpha)$ interpolates:
\begin{equation}
  \mathbf{l} = \alpha\,\mathbf{l}_{\mathrm{rule}} +
               (1{-}\alpha)\,\mathbf{l}_{\mathrm{neural}}.
\end{equation}
Optional $\alpha$-regularisation pushes the gate toward full interpretability:
\begin{equation}
  \mathcal{L}_\alpha = \mathcal{L}_{\mathrm{CE}} + \lambda_\alpha(1{-}\alpha)^2.
  \label{eq:alpha_reg}
\end{equation}

\subsubsection{OOD Scoring in Rule-Activation Space}

At inference, binary rule activations $\mathbf{a}^{\mathrm{hard}}(\mathbf{x})$
form an $R$-dimensional embedding.
We fit the Mahalanobis detector (Eq.~\ref{eq:mahal}) in this space using
known-class training activations.
Novel attacks are expected to activate different rule subsets than known
classes, producing anomalous activation patterns detectable by distance.

\subsection{Baselines and Evaluation}

An unconstrained two-hidden-layer MLP (256, 128 units, ReLU, dropout 0.3)
serves as the accuracy ceiling.
A Decision Tree (DT; depth search over $\{5,10,15,\infty\}$, selected by
validation F1-score (F1)) and a Random Forest (RF;~\cite{breiman2001random}
100 trees, depth 10) represent classical rule-based baselines.
DT and RF OOD scoring uses both maximum-confidence ($1{-}\max_c p_c$) and
Mahalanobis distance in the raw feature space.

\textbf{Metrics:} weighted F1 on known-class test traffic;
AUROC (area under the receiver operating characteristic curve) for OOD
detection (unknown samples as positive); true-positive rate at 5\% false
positive rate (TPR@5\%FPR); rule crispness; and gate value $\alpha$.

\textbf{Training:} Adam optimiser~\cite{kingma2014adam}
($\mathrm{lr}{=}10^{-3}$, batch size 512, up to 60 epochs, patience 12
on validation F1).
Five independent random seeds (0--4) for NeSy-NIDS and CBM ($\gamma{=}0.5$);
single run for all other models.

% ─── 4. Results ───────────────────────────────────────────────────────────────
\section{Results and Discussion}

\subsection{Classification Accuracy}

Table~\ref{tab:main} summarises all results.
Both NeSy-NIDS and CBM ($\gamma{=}0.5$) \textbf{match the unconstrained MLP}
on both datasets.
The F1 gap between NeSy and MLP is +0.0013 on CTU-IoT-23 and +0.0015 on
CIC-IoT-2023---well within one standard deviation.
This replicates the zero-accuracy-cost claim of Koh et al.~\cite{koh2020concept}
in the IoT NIDS domain and demonstrates that the constraints imposed by
interpretable architectures carry no performance penalty.

The Rule-only ablation degrades substantially on CTU-IoT-23 (F1\,=\,0.5612),
confirming that the neural fallback handles traffic patterns outside the
eight rule templates, and that the $\alpha$-gate correctly delegates these
cases. The fallback is less critical on CIC-IoT-2023 (F1\,=\,0.7561) where
the tighter CIC rules cover more of the class structure.

DT and RF achieve comparable F1 (DT: 0.9350/0.8116; RF: 0.9351/0.8045) but
at severe cost to OOD detection, as detailed next.

\begin{table*}[!t]
\renewcommand{\arraystretch}{1.15}
\caption{Results on CTU-IoT-23 and CIC-IoT-2023. Mean $\pm$ std over 5 seeds
where available; single run otherwise. AUROC uses Mahalanobis distance in
the respective representation space (rule/concept activations for NeSy/CBM;
raw features for DT/RF~$\dagger$).}
\label{tab:main}
\centering
\small
\begin{tabular}{lcccccc}
\toprule
& \multicolumn{3}{c}{\textbf{CTU-IoT-23}} &
  \multicolumn{3}{c}{\textbf{CIC-IoT-2023}} \\
\cmidrule(lr){2-4}\cmidrule(lr){5-7}
\textbf{Model} & \textbf{F1} & \textbf{AUROC} & \textbf{TPR@5\%} &
                 \textbf{F1} & \textbf{AUROC} & \textbf{TPR@5\%} \\
\midrule
MLP Baseline
  & 0.9323 & 0.913 & 0.596
  & 0.8188 & 0.858 & 0.715 \\
JointCBM ($\gamma{=}0$)
  & 0.9322 & 0.882 & 0.426
  & 0.8168 & 0.560 & 0.043 \\
JointCBM ($\gamma{=}0.5$)
  & 0.9325 & $0.918\pm0.019$ & 0.687
  & 0.8154 & 0.591 & 0.102 \\
\textbf{NeSy-NIDS}
  & \textbf{0.9336} & $\mathbf{0.909\pm0.013}$ & 0.553
  & \textbf{0.8203} & $\mathbf{0.614\pm0.011}$ & 0.227 \\
NeSy + $\alpha$-reg
  & 0.9336 & $0.895\pm0.023$ & 0.505
  & 0.8205 & $0.634\pm0.021$ & 0.236 \\
Rule-only
  & 0.5612 & 0.863 & 0.459
  & 0.7561 & 0.693 & 0.270 \\
DT$^\dagger$
  & 0.9350 & 0.335 & 0.037
  & 0.8116 & 0.809 & 0.560 \\
RF$^\dagger$
  & 0.9351 & 0.627 & 0.411
  & 0.8045 & 0.694 & 0.140 \\
\bottomrule
\end{tabular}
\end{table*}

\subsection{Open-Set Detection}

Fig.~\ref{fig:auroc_vs_f1} plots OOD AUROC against F1, revealing a clear
separation between model families that depends on dataset structure.

\begin{figure}[!t]
  \centering
  \includegraphics[width=\columnwidth]{figures/auroc_vs_f1.pdf}
  \caption{OOD AUROC vs.\ known-class weighted F1 for all models on
  CTU-IoT-23 (left) and CIC-IoT-2023 (right). Error bars indicate
  $\pm$1 std over 5 seeds.}
  \label{fig:auroc_vs_f1}
\end{figure}

\textbf{CTU-IoT-23.}
Rule-space and concept-space Mahalanobis scoring are both highly effective.
NeSy-NIDS achieves AUROC\,=\,0.909 (std\,=\,0.013)---the lowest variance
of any model---while CBM ($\gamma{=}0.5$) reaches 0.918 (std\,=\,0.019).
Both substantially outperform the DT (0.335) and RF (0.627), which use
Mahalanobis on the raw 37-dimensional feature space.
The DT AUROC of 0.335 is below random (0.5), indicating that greedy
classification splits produce representations that are \emph{harder} to
use for OOD separation than either the rule or concept space.
This is the primary positive result: enforcing interpretability does not
reduce OOD detection---it improves it relative to classical rule-based models.

\textbf{CIC-IoT-2023.}
All bottleneck models show degraded OOD detection relative to the unconstrained
MLP (0.858).
NeSy-NIDS achieves 0.614 (std\,=\,0.011) and CBM ($\gamma{=}0.5$) achieves 0.591.
The DT, however, achieves 0.809 using Mahalanobis on the raw 39-dimensional
feature space.
This is the \textbf{compression-structural OOD failure}: reducing 39 features
to 8 rule or concept activations discards information necessary to
distinguish CIC's 29 unknown attack families from known traffic.
Both bottleneck architectures independently replicate this failure (NeSy: 0.614,
CBM: 0.591), confirming it is a property of the compression, not an
implementation artefact.
As argued by Fort et al.~\cite{fort2021exploring}, representation quality
governs Mahalanobis OOD performance; here the quality is bottlenecked by
the limited concept/rule vocabulary, not the neural backbone.

The CIC result aligns with the open-world findings of Han et
al.~\cite{han2023owad}: not every representation space is equally suited
for detecting distributional novelty. The 29 unknown CIC attack classes
activate similar rule subsets to known attacks, rendering the 8-dimensional
rule space non-discriminative for OOD purposes.

\subsection{Rule Crispness and Class Selectivity}

After $k$-annealing to $k{=}10$ and STE binarisation, \textbf{100\% of rule
activations on both datasets fall in $[0,0.1)\cup(0.9,1]$}, confirming
exact Boolean behaviour at inference.

Figs.~\ref{fig:sel_ctu} and~\ref{fig:sel_cic} show the mean rule activation
per class (class selectivity heatmaps).
On CTU-IoT-23, \texttt{incomplete\_handshake} fires at 0.859 for C\&C-HeartBeat
(SYN-heavy flows with few ACKs), \texttt{asymmetric\_flood} at 0.796 for
DDoS, and \texttt{cc\_heartbeat\_beaconing} at 0.772 for C\&C-HeartBeat.
These values are semantically consistent with the rule definitions, validating
that the learned thresholds converge to meaningful boundaries.

On CIC-IoT-2023, selectivity is extreme: \texttt{syn\_flood\_ddos} activates
at 0.996 for DDoS-SYN and \texttt{udp\_mirai\_flood} at 0.994 for Mirai
traffic.
High selectivity explains both the strong classification performance and the
OOD detection failure: rules that fire with near-certainty for specific
known classes provide strong classification signal but little distributional
spread for Mahalanobis separation.

\begin{figure}[!t]
  \centering
  \begin{subfigure}[t]{0.49\columnwidth}
    \includegraphics[width=\textwidth]{figures/rule_selectivity_ctu.pdf}
    \caption{CTU-IoT-23}
    \label{fig:sel_ctu}
  \end{subfigure}\hfill
  \begin{subfigure}[t]{0.49\columnwidth}
    \includegraphics[width=\textwidth]{figures/rule_selectivity_cic.pdf}
    \caption{CIC-IoT-2023}
    \label{fig:sel_cic}
  \end{subfigure}
  \caption{Rule class selectivity heatmaps. Cell values are mean rule
  activation per class at $k{=}10$ (post-STE). High values (dark)
  indicate selective activation for a class.}
  \label{fig:sel}
\end{figure}

\subsection{Threshold Drift}

Fig.~\ref{fig:drift} shows learned threshold drift (learned $-$ initial)
per rule condition.
On CTU-IoT-23, \texttt{cc\_heartbeat\_beaconing} inter-arrival time (IAT)
bounds shift $-0.2$ to $-0.4$ (the model learns a wider beaconing window),
\texttt{asymmetric\_flood} response-packet threshold shifts $+0.17$ (admits
more borderline asymmetric flows), and \texttt{incomplete\_handshake} SYN
threshold shifts $+0.19$ (tighter SYN requirement).
All drifts are semantically consistent with what domain knowledge would
predict from the data distribution, representing principled refinement of
expert-defined rules rather than arbitrary drift.

\begin{figure}[!t]
  \centering
  \begin{subfigure}[t]{0.49\columnwidth}
    \includegraphics[width=\textwidth]{figures/threshold_drift_ctu.pdf}
    \caption{CTU-IoT-23}
    \label{fig:drift_ctu}
  \end{subfigure}\hfill
  \begin{subfigure}[t]{0.49\columnwidth}
    \includegraphics[width=\textwidth]{figures/threshold_drift_cic.pdf}
    \caption{CIC-IoT-2023}
    \label{fig:drift_cic}
  \end{subfigure}
  \caption{Learned threshold drift (learned $-$ initial) per rule condition.
  Red (positive): threshold increased. Blue (negative): threshold decreased.}
  \label{fig:drift}
\end{figure}

\subsection{The $\alpha$-Gate and Interpretability Trade-Off}

Without regularisation, NeSy-NIDS converges to $\alpha{=}0.37$ on CTU-IoT-23
and $\alpha{=}0.27$ on CIC-IoT-2023.
With $\lambda_\alpha{=}0.1$, the gate rises to $\alpha{=}0.68$ and 0.85
respectively at zero F1 cost.
On CTU-IoT-23, higher $\alpha$ incurs a minor OOD cost (AUROC: 0.909 $\to$
0.895), consistent with the rule head taking over from a neural head whose
embedding was more OOD-compact.
On CIC-IoT-2023, OOD AUROC marginally improves under $\alpha$-reg
(0.614 $\to$ 0.634), since the rule-dominant embedding is tighter.

The $\alpha$-gate provides a \emph{deployable Pareto front}: an operator
can select any interpretability-OOD trade-off point at inference time by
adjusting $\lambda_\alpha$ during a brief fine-tuning step, without full
retraining. This directly addresses the operator-control desideratum
identified by Arp et al.~\cite{arp2022dos}.

\subsection{Concept Leakage in CBMs}

JointCBM ($\gamma{=}0$) achieves F1\,=\,0.9322 on CTU-IoT-23, matching
the MLP. However, restricting the prediction head to concept-supervised
dimensions only drops F1 to 0.7810---a 15-point gap indicating that the
joint model encodes class-discriminative information in the concept layer
beyond the 10 specified concepts.
With $\gamma{=}0.5$, this gap closes substantially (F1\,=\,0.9325) while
OOD AUROC improves from 0.882 to 0.918.
This confirms that concept leakage and OOD detection quality are
negatively correlated: when concepts are forced to be semantically
faithful, the concept space becomes more compact and Mahalanobis
distances are more discriminative.
This is, to our knowledge, the first explicit quantification of concept
leakage and its effect on OOD detection in the security domain,
complementing the post-hoc explanation fidelity analysis of xNIDS~\cite{wei2023xnids}
and CADE~\cite{yang2021cade}.

\subsection{Discussion}

\textbf{Dataset structure governs OOD detectability.}
CTU-IoT-23 has four sharply delineated known classes with
traffic profiles that map cleanly onto eight expert rules.
CIC-IoT-2023's 29 unknown attack families include many that share
high-level characteristics (high packet rate, short connections) with
known attacks, causing rule-space overlap.
This aligns with the observation of Fort et al.~\cite{fort2021exploring}
that representation quality---here the coverage of the rule vocabulary---
dominates OOD performance.

\textbf{Inherent vs.\ post-hoc interpretability.}
xNIDS~\cite{wei2023xnids} achieves high-fidelity post-hoc explanations
from a black-box model. Our work demonstrates that inherently interpretable
models match black-box accuracy, while additionally supporting OOD
detection in the interpretable space and enabling test-time concept
intervention (CBM) or threshold audit (NeSy).
The two approaches are complementary: a deployment prioritising operator
intervention would prefer CBM; one requiring forensic threshold inspection
would prefer NeSy-NIDS.

\textbf{Comparison to open-world NIDS.}
OWAD~\cite{han2023owad} targets normality shift detection in deployed
anomaly detectors via hypothesis testing over time.
Our setting is distinct: we target \emph{unknown class detection} at
single-inference granularity, in inherently interpretable models.
Together, the two approaches cover the lifecycle of a deployed NIDS:
our method flags unknown-class traffic at inference; OWAD detects when
the distribution of known traffic has shifted and retraining is warranted.

% ─── 5. Conclusion ────────────────────────────────────────────────────────────
\section{Conclusion}

We have presented a joint empirical evaluation of CBMs and NeSy-NIDS for
interpretable open-set IoT NIDS on CTU-IoT-23 and CIC-IoT-2023.
Both architectures match unconstrained MLP classification accuracy while
providing human-auditable intermediate representations.
On CTU-IoT-23, both achieve competitive OOD AUROC (up to 0.918),
substantially outperforming classical rule-based models.
On CIC-IoT-2023, a structural compression effect limits OOD detection
for both---a finding we trace to the mismatch between an 8-rule vocabulary
and 29 diverse unknown attack families.

NeSy-NIDS additionally contributes 100\% crisp binary rule activations,
semantically consistent threshold drift, and a controllable $\alpha$-gate.
CBMs contribute concept leakage quantification and elimination, with
$\gamma$-regularisation simultaneously improving concept fidelity and OOD
detection quality.

Future work should explore expanding rule and concept vocabularies through
data-driven template discovery, combining NeSy and CBM bottlenecks in a
hierarchical architecture, and investigating how the findings generalise to
enterprise-scale traffic where the ratio of known to unknown attack classes
is less favourable.

\bibliographystyle{IEEEtran}
\bibliography{refs}

\end{document}
